\section*{Chapter \ref{ch:g11:lithosphere-asthenosphere} --- Lithosphere and Asthenosphere}

\begin{solution}{exo:g11:lithosphere-asthenosphere:1}
The \omterm{def:g11:lithosphere-asthenosphere:lithosphere}{lithosphere} is the rigid outer shell of the Earth, made of the crust and
the uppermost mantle; the \omterm{def:g11:lithosphere-asthenosphere:lithosphere}{asthenosphere} is the hotter, weaker mantle beneath
it, which flows slowly over long times. The oceanic crust, the continental
crust and the uppermost mantle belong to the \omterm{def:g11:lithosphere-asthenosphere:lithosphere}{lithosphere}; the outer core does
not.
\end{solution}

\begin{solution}{exo:g11:lithosphere-asthenosphere:2}
The iceberg sinks until it pushes aside its own weight of water: the
fraction under water is $910/1025 \approx 0.888$, about \qty{89}{\percent}.
The fraction above is about \qty{11.2}{\percent}, so the total height is
$20/0.112 \approx \qty{180}{m}$, of which about \qty{160}{m} lie under
water.
\end{solution}

\begin{solution}{exo:g11:lithosphere-asthenosphere:3}
$15 + 26 \times 3 = \qty{93}{\celsius}$. At \qty{100}{km} the same gradient
would give more than \qty{2600}{\celsius}, far above the melting temperature
of mantle rock; but \omterm{def:g11:seismic-structure:body-waves}{S waves} cross the upper mantle, which is therefore solid.
The gradient must become much smaller with depth.
\end{solution}

\begin{solution}{exo:g11:lithosphere-asthenosphere:4}
The band from 0 to \qty{500}{m} above sea level, with about
\qty{16}{\percent} of the surface, and the band from 4000 to \qty{4500}{m}
below sea level, with about \qty{11}{\percent} (the band from 4500 to
\qty{5000}{m} below is almost as large).
\end{solution}

\begin{solution}{exo:g11:lithosphere-asthenosphere:5}
At $P = 0$: \qty{1121}{\celsius}. At \qty{3}{GPa}: $-5.104 \times 9 +
132.899 \times 3 + 1120.661 \approx -45.9 + 398.7 + 1120.7 \approx
\qty{1473}{\celsius}$.
\end{solution}

\begin{solution}{exo:g11:lithosphere-asthenosphere:6}
About \qty{730}{\celsius} at \qty{50}{km} (\omterm{def:g11:lithosphere-asthenosphere:lithosphere}{lithosphere}), about
\qty{1370}{\celsius} at \qty{150}{km} and \qty{1410}{\celsius} at
\qty{250}{km} (\omterm{def:g11:lithosphere-asthenosphere:lithosphere}{asthenosphere}). The temperature rises most between 50 and
\qty{150}{km}, by more than \qty{600}{\celsius}, because most of that
interval lies in the \omterm{def:g11:lithosphere-asthenosphere:lithosphere}{lithosphere}; between 150 and \qty{250}{km} it rises by
only about \qty{45}{\celsius}.
\end{solution}

\begin{solution}{exo:g11:lithosphere-asthenosphere:7}
Root: $3 \times 2830/550 \approx \qty{15.4}{km}$. Extra thickness:
$3 \times 3380/550 \approx \qty{18.4}{km}$, that is the \qty{3}{km} of the
range plus its root.
\end{solution}

\begin{solution}{exo:g11:lithosphere-asthenosphere:8}
The crust is thicker by $55 - 41.1 = \qty{13.9}{km}$, and $\Delta H = h
\times 3380/550$, so $h = 13.9 \times 550/3380 \approx \qty{2.3}{km}$.
\end{solution}

\begin{solution}{exo:g11:lithosphere-asthenosphere:9}
Decrease: $(4.470 - 4.419)/4.470 \approx \qty{1.1}{\percent}$. Jump:
$(4.644 - 4.419)/4.419 \approx \qty{5.1}{\percent}$. With depth the rock is
more and more squeezed, which makes seismic waves faster; a decrease means
that something else wins: the rock is hot and close to melting, and weaker.
\end{solution}

\begin{solution}{exo:g11:lithosphere-asthenosphere:10}
(1) $(1345 - 610)/60 \approx \qty{12}{\celsius}$ per kilometre. (2) That is
about half the gradient at the surface. Part of the heat is made inside the
upper crust; that heat must also cross the top layers to escape, which makes
the gradient steeper near the surface. Below the crust only the heat coming
from deeper down crosses the rock, and the gradient is smaller.
\end{solution}

\begin{solution}{exo:g11:lithosphere-asthenosphere:11}
At \qty{200}{km}: $1300 + 90 = \qty{1390}{\celsius}$; at \qty{400}{km}:
$1300 + 180 = \qty{1480}{\celsius}$. The rise is \qty{90}{\celsius} over
\qty{200}{km}. Through the \omterm{def:g11:lithosphere-asthenosphere:lithosphere}{lithosphere} the temperature rises by
$1345 - 15 = \qty{1330}{\celsius}$ over only \qty{100}{km}: about thirty
times more per kilometre.
\end{solution}

\begin{solution}{exo:g11:lithosphere-asthenosphere:12}
(1) $118 + 12.3\,z = 1300 + 0.45\,z$ gives $z = 1182/11.85 \approx
\qty{100}{km}$. (2) $118 + 21.5\,z = 1300 + 0.45\,z$ gives $z = 1182/21.05
\approx \qty{56}{km}$. (3) The hotter west has a much thinner \omterm{def:g11:lithosphere-asthenosphere:lithosphere}{lithosphere},
as measured (about \qty{40}{km} at the thinnest); the very thick \omterm{def:g11:lithosphere-asthenosphere:lithosphere}{lithosphere}
of the middle of the continent needs an even smaller gradient.
\end{solution}

\begin{solution}{exo:g11:lithosphere-asthenosphere:13}
(1) By the proposition, $\Delta H = h\,\rho_m/(\rho_m - \rho_c)$, so $h =
\Delta H\,(\rho_m - \rho_c)/\rho_m$. (2) $550/3380 \approx 0.16$: about one
sixth of any extra crust shows as height, five sixths as root. (3) For the
iceberg, the fraction above the water is $(1025 - 910)/1025 \approx 0.11$:
the same rule, (density of the fluid $-$ density of the floating body)/density
of the fluid.
\end{solution}

\begin{solution}{exo:g11:lithosphere-asthenosphere:14}
(1) $1489 - 1345 = \qty{144}{\celsius}$. (2) A little melt gathers in thin
films between the grains of the rock. The grains are then no longer firmly
locked together: the rock resists shear less, so \omterm{def:g11:seismic-structure:body-waves}{S waves} slow down, and it
can deform and flow far more easily.
\end{solution}

\begin{solution}{exo:g11:lithosphere-asthenosphere:15}
(1) No. Planing the land would make the continental crust a little thinner,
but it would still be thick and light, and would still float much higher than
the thin, dense crust of the ocean floor; the oceans would still fill the
basins between the continents. (2) There are two kinds of crust, of different
thickness and density: a thick, light continental crust and a thin, dense
oceanic crust.
\end{solution}

\begin{solution}{pb:g11:lithosphere-asthenosphere:1}
\textbf{1.} $5.01 - 0.67 = \qty{4.34}{km}$.

\textbf{2.} $(3380 - 2830)/3380 \approx \qty{16}{\percent}$.

\textbf{3.} The fraction $2830/3380 \approx 0.84$ would lie below; the
iceberg sinks a little deeper, $910/1025 \approx 0.89$, because ice is
closer in density to water than crust is to mantle.

\textbf{4.} For the seconds of a seismic wave, or the time of an
earthquake, the mantle is a solid; only over very long times does the
\omterm{def:g11:lithosphere-asthenosphere:lithosphere}{asthenosphere} flow and let the crust float.

\textbf{5.} Below the deepest \omterm{def:g11:seismic-structure:moho}{Moho}, both columns are made of the same mantle
rock at the same depths, and they weigh the same there; all the differences
lie above that depth and must be included.

\textbf{6.} The average continent: crust from \qty{0.67}{km} above to
\qty{40.4}{km} below sea level, then mantle down to \qty{70}{km}. The
plateau: crust from \qty{5.01}{km} above sea level down to a \omterm{def:g11:seismic-structure:moho}{Moho} deeper by
the root, then mantle down to \qty{70}{km}.

\textbf{7.} The plateau column has $h + r$ more crust, the continent column
has mantle over the thickness $r$ instead: $\rho_c\,(h + r) = \rho_m\,r$,
so $r = h\,\rho_c/(\rho_m - \rho_c)$ and $\Delta H = h + r =
h\,\rho_m/(\rho_m - \rho_c)$.

\textbf{8.} $3380/550 \approx 6.15$.

\textbf{9.} $\Delta H = 4.34 \times 6.15 \approx \qty{26.7}{km}$; the root is
$26.7 - 4.34 \approx \qty{22.3}{km}$.

\textbf{10.} The average continent's \omterm{def:g11:seismic-structure:moho}{Moho} lies $41.1 - 0.67 \approx
\qty{40.4}{km}$ below sea level; the plateau's lies $40.4 + 22.3 \approx
\qty{62.8}{km}$ below sea level.

\textbf{11.} $41.1 + 26.7 \approx \qty{67.8}{km}$.

\textbf{12.} One kilometre of rock above a square metre has a mass of
$\rho \times \qty{1000}{m^3}$: \num{2830} tonnes of crust or \num{3380}
tonnes of mantle. Continent: $41.1 \times 2830 + (70 - 40.43) \times 3380
\approx \num{116300} + \num{99900} \approx \num{216300}$ tonnes. Plateau:
$67.77 \times 2830 + (70 - 62.76) \times 3380 \approx \num{191800} +
\num{24500} \approx \num{216300}$ tonnes. The two are equal.

\textbf{13.} $41.1 - 0.67 \times 6.15 \approx \qty{37.0}{km}$.

\textbf{14.} $67.8/41.1 \approx 1.65$ times.

\textbf{15.} $3300/470 \approx 7.02$, $\Delta H \approx 4.34 \times 7.02
\approx \qty{30.5}{km}$, and a crust of about \qty{71.6}{km}. A change of
about \qty{2}{\percent} in the mantle's density changes the thickness by
about \qty{4}{km}: the result depends noticeably on the densities, but stays
close to \qty{70}{km}.

\textbf{16.} The model gives about \qty{68}{km}, the seismic waves about
\qty{70}{km}: they agree within a few per cent.

\textbf{17.} $h = 3.23 - 0.67 = \qty{2.56}{km}$, $\Delta H = 2.56 \times
6.15 \approx \qty{15.7}{km}$, and a crust of about \qty{56.8}{km}, some
\qty{7}{km} more than measured.

\textbf{18.} Thinner. A column whose upper part is lighter than average
weighs less for the same thickness, so it floats higher: a lighter crust
needs a smaller root to stand at a given height. The real crust can
therefore be thinner than the model's, as measured.

\textbf{19.} Over long times the \omterm{def:g11:lithosphere-asthenosphere:lithosphere}{asthenosphere} flows: it moves aside under
a heavy column and comes back under a light one, until the columns balance.

\textbf{20.} \textbf{The crust under Tibet is about \qty{70}{km} thick}, much
more than the \qty{41}{km} of an average continent: the plateau stands five
kilometres high because its thick, light crust floats on the denser mantle,
with a root of more than \qty{20}{km}.
\end{solution}
